\vspace{0.5em}\noindent\textbf{Mục tiêu Tuần 3:}

\begin{itemize}
\tightlist
\item
  Tìm hiểu và nghiên cứu về bài toán trích xuất thực thể y tế tiếng
  Việt.
\item
  Xử lý dữ liệu văn bản y tế đầu vào.
\item
  Tìm hiểu các kiến trúc học sâu cơ bản và xây dựng mô hình AI nhận diện
  thực thể y tế (ViMQ).
\end{itemize}

\vspace{0.5em}\noindent\textbf{Các công việc triển khai trong tuần:}

\begingroup
\small
\setlength{\tabcolsep}{3pt}
\renewcommand{\arraystretch}{1.15}
\begin{longtable}[]{@{}
  >{\raggedright\arraybackslash}p{(\columnwidth - 8\tabcolsep) * \real{0.2000}}
  >{\raggedright\arraybackslash}p{(\columnwidth - 8\tabcolsep) * \real{0.2000}}
  >{\raggedright\arraybackslash}p{(\columnwidth - 8\tabcolsep) * \real{0.2000}}
  >{\raggedright\arraybackslash}p{(\columnwidth - 8\tabcolsep) * \real{0.2000}}
  >{\raggedright\arraybackslash}p{(\columnwidth - 8\tabcolsep) * \real{0.2000}}@{}}
\toprule\noalign{}
\begin{minipage}[b]{\linewidth}\raggedright
Thứ
\end{minipage} & \begin{minipage}[b]{\linewidth}\raggedright
Công việc
\end{minipage} & \begin{minipage}[b]{\linewidth}\raggedright
Ngày bắt đầu
\end{minipage} & \begin{minipage}[b]{\linewidth}\raggedright
Ngày hoàn thành
\end{minipage} & \begin{minipage}[b]{\linewidth}\raggedright
Nguồn tài liệu
\end{minipage} \\
\midrule\noalign{}
\endhead
\bottomrule\noalign{}
\endlastfoot
\textbf{2} & Phân tích đặc thù văn bản y tế và làm quen với bộ dữ liệu y
khoa tiếng Việt. & 15/06/2026 & 15/06/2026 &
\href{https://github.com/tadeephuy/ViMQ}{ViMQ Repository} \\
\textbf{3} & Viết mã tiền xử lý dữ liệu để chuyển đổi văn bản thành các
con số cho AI hiểu (Tokenization \& Encoding). & 16/06/2026 & 16/06/2026
& \href{https://github.com/tadeephuy/ViMQ}{ViMQ Repository} \\
\textbf{4} & Chuẩn bị và biến đổi dữ liệu nhãn thành định dạng chuẩn để
huấn luyện mô hình. & 17/06/2026 & 17/06/2026 &
\href{https://github.com/tadeephuy/ViMQ}{ViMQ Repository} \\
\textbf{5} & Lên ý tưởng và bắt đầu lắp ghép các thành phần mạng nơ-ron
cơ bản cho mô hình ViMQ. & 18/06/2026 & 19/06/2026 &
\href{https://github.com/tadeephuy/ViMQ}{ViMQ Repository} \\
\textbf{6} & Hoàn thiện cấu trúc mô hình để dự đoán được loại bệnh,
triệu chứng từ câu văn. & 20/06/2026 & 21/06/2026 &
\href{https://github.com/tadeephuy/ViMQ}{ViMQ Repository} \\
\end{longtable}
\endgroup

\vspace{0.5em}\noindent\textbf{Kết quả đạt được:}

\begin{itemize}
\tightlist
\item
  Hoàn thành xây dựng pipeline tiền xử lý dữ liệu NER span-based chuyên
  sâu cho ngôn ngữ tiếng Việt.
\item
  Lập trình thành công kiến trúc mô hình ViMQ hiện đại, kết hợp sức mạnh
  của PhoBERT, BiLSTM và Biaffine Classifier.
\item
  Đặt nền móng vững chắc cho khâu phân tích câu hỏi chuyên môn y tế
  (Query Analyzer) trong hệ sinh thái Smart Healthcare Platform.
\end{itemize}
